\section{Tối ưu hóa truy xuất dữ liệu}

Dựa trên pipeline xử lý truy vấn của hệ quản trị CSDL, việc viết câu lệnh SQL là chưa đủ. Hệ thống cần tối ưu hóa cây truy vấn (Query Tree) thông qua các luật biến đổi tương đương để giảm thiểu chi phí bộ nhớ và I/O trước khi lập kế hoạch thực thi (Execution Plan).

\subsection{Bài toán truy vấn}
\textbf{Yêu cầu:} Lấy ra Mã giao dịch (\texttt{TransactionID}) và Số điểm (\texttt{Amount}) của các giao dịch cộng điểm lớn hơn 500, thuộc về các Chiến dịch (\texttt{CAMPAIGN}) đang ở trạng thái hoạt động (\texttt{'Active'}).

\textbf{Biểu thức đại số quan hệ:}
\[
\begin{aligned}
T' &= \sigma_{Amount > 500 \,\wedge\, TransactionType = \text{'Earn'}}(T),\\
C' &= \sigma_{Status = \text{'Active'}}(C),\\
Q &= \pi_{TransactionID,Amount}(T' \bowtie_{T'.CampaignID=C'.CampaignID} C').
\end{aligned}
\]
Trong đó $T$ là CREDIT\_TRANSACTION, $C$ là CAMPAIGN. Phép nối dùng điều kiện bằng CampaignID, không phải nối tự nhiên trên tất cả thuộc tính cùng tên.

\subsection{Cây truy vấn trước và sau khi tối ưu}

Theo luật biến đổi tương đương, phép chọn (Selection - $\sigma$) nếu được đẩy xuống trước phép kết nối (Join - $\bowtie$) sẽ giúp giảm thiểu đáng kể kích thước dữ liệu trung gian.

\begin{figure}[H]
    \centering
    \resizebox{\textwidth}{!}{ % Ép toàn bộ sơ đồ tự động co lại vừa khít 100% chiều ngang trang giấy
    \begin{tikzpicture}[
        box/.style={rectangle, draw=black, rounded corners, inner sep=12pt, text centered, font=\normalsize},
        leaf/.style={rectangle, draw=black, inner sep=10pt, text centered, font=\normalsize\bfseries}
    ]
        % ================= CÂY CHƯA TỐI ƯU =================
        \node[font=\Large\bfseries] at (0, 8.5) {Chưa tối ưu};
        
        \node[box] (proj1) at (0, 7) {$\pi_{TransactionID, Amount}$};
        \node[box] (sel1) at (0, 5) {$\sigma_{Amount > 500 \wedge TransactionType = \text{'Earn'} \wedge Status = \text{'Active'}}$};
        \node[box] (join1) at (0, 3) {$\bowtie_{CampaignID}$};
        
        \node[leaf] (trans1) at (-4.5, 0) {CREDIT\_TRANSACTION};
        \node[leaf] (camp1) at (4.5, 0) {CAMPAIGN};
        
        \draw[->, thick] (sel1) -- (proj1);
        \draw[->, thick] (join1) -- (sel1);
        \draw[->, thick] (trans1) -- (join1);
        \draw[->, thick] (camp1) -- (join1);

        % ================= CÂY ĐÃ TỐI ƯU =================
        \node[font=\Large\bfseries] at (18, 8.5) {Đã đẩy phép chọn};
        
        \node[box] (proj2) at (18, 7) {$\pi_{TransactionID, Amount}$};
        \node[box] (join2) at (18, 5) {$\bowtie_{CampaignID}$};
        
        \node[box] (sel2_trans) at (13, 2.5) {$\sigma_{Amount > 500 \wedge TransactionType = \text{'Earn'}}$};
        \node[box] (sel2_camp) at (23, 2.5) {$\sigma_{Status = \text{'Active'}}$};
        
        \node[leaf] (trans2) at (13, 0) {CREDIT\_TRANSACTION};
        \node[leaf] (camp2) at (23, 0) {CAMPAIGN};
        
        \draw[->, thick] (join2) -- (proj2);
        \draw[->, thick] (sel2_trans) -- (join2);
        \draw[->, thick] (sel2_camp) -- (join2);
        \draw[->, thick] (trans2) -- (sel2_trans);
        \draw[->, thick] (camp2) -- (sel2_camp);
        
        % Đường phân cách dọc ở giữa
        \draw[dashed, thick, gray] (9, -1) -- (9, 9.5);
    \end{tikzpicture}
    }
    \caption{Cây truy vấn trước và sau khi tối ưu hóa bằng luật phân rã}
\end{figure}

\subsection{Câu lệnh SQL tương ứng}
\begin{lstlisting}[style=sqlstyle]
SELECT ct.TransactionID, ct.Amount
FROM CREDIT_TRANSACTION AS ct
JOIN CAMPAIGN AS c ON c.CampaignID = ct.CampaignID
WHERE ct.Amount > 500
  AND ct.TransactionType = 'Earn'
  AND c.Status = 'Active';
\end{lstlisting}

\subsection{Đánh giá thuật toán kết nối}
Đẩy phép chọn xuống có thể giảm số dòng trung gian; mức giảm thực tế phụ thuộc phân bố dữ liệu. SQL Server chọn Nested Loops, Hash hoặc Merge Join theo statistics, chỉ mục và chi phí ước lượng. Nested Loops thường phù hợp khi đầu vào ngoài nhỏ và phía trong có đường tra cứu hiệu quả; không thể kết luận chỉ từ việc CAMPAIGN nhỏ.

Đồ án chưa có execution plan hoặc số đo I/O, thời gian trước/sau. Vì vậy đây là phân tích biến đổi logic, không phải bằng chứng một thuật toán luôn được chọn hay truy vấn đã đạt mục tiêu hiệu năng.
